\part*{Part II --- Four containers, four logics}
\addcontentsline{toc}{part}{Part II --- Four containers, four logics}

\bigskip
\noindent
\emph{This part works the collection dial. It shows first that the algebraic
theory presenting the container \emph{is} the substructural hierarchy of the
induced propositional logic; then that a second, finer axis is needed --- which
lattice operations the container makes primitive --- and that three settings of
that axis give intuitionistic logic, classical linear logic, and the logic of
decentralised consensus. Linear logic is worked in detail, because it is the
setting at which the framework's claim to be generating rather than modelling is
most testable: its two conjunctions, its exponential and its involutive negation
each turn out to name a feature of the container. The part closes by supplying
the structural and behavioural components for the coalition-logic row, which
\cite{Zeus} deferred to a companion. Sources: \cite{Zeus} throughout;
\cite{Girard87,Yetter90,Rosenthal90} for the linear row; \cite{Gabbay25,
GabbaySemiframes,QuorumNote,ConsensusTypes} for the semitopological row.}
\bigskip

\section{The substructural dial}
\label{sec:dial}

The collection type $\Coll$ is a monad, and the algebraic theory presenting it is
not incidental: \emph{its equations are the structural rules of the propositional
logic the collection component induces}. Write the binary aggregation as $\ast$
with unit $\varepsilon$. Then
\[
\begin{array}{lll}
  \text{commutativity of } \ast & \longleftrightarrow & \text{exchange},\\[2pt]
  \text{idempotence of } \ast    & \longleftrightarrow & \text{contraction},\\[2pt]
  \text{unit and absorption laws}    & \longleftrightarrow & \text{weakening}.
\end{array}
\]

Reading the standard containers through this correspondence:

\begin{itemize}[leftmargin=1.4em]
\item $\Pfin$, finite powerset --- a commutative idempotent monoid, i.e.\ a
join-semilattice --- gives exchange and contraction, and with the unit, weakening:
the full structural regime, hence Heyting or Boolean propositional logic.
\item $\Mset$, finite multisets --- the free commutative monoid, commutative but
not idempotent --- gives exchange without contraction: the multiplicative-linear
regime, with $\otimes$.
\item $\List$ --- the free monoid, neither commutative nor idempotent --- gives
neither: a non-commutative, Lambek-style regime \cite{Lambek58}.
\item $\Dist$, distributions --- commutative, with a convex structure in place of
idempotence --- gives the probabilistic regime.
\end{itemize}

\begin{principle}[Collection is substructure]
\label{prin:dial}
The setting of the collection dial is the pair of Booleans
$(\text{commutative?},\ \text{idempotent?})$ of $\Coll$, and that pair is exactly
the presence or absence of exchange and contraction in the induced logic.
\end{principle}

Two coordinates, and everything in this part is bookkeeping on them. It is worth
noticing already that the same two coordinates are the ones on which the no-go
theorems for distributive laws between monads turn \cite{ZwartMarsden19}: the
obstruction to composing two containers and the price of packaging a container's
own sub-collections back into itself are one datum read twice.

\begin{remark}[Commutative only]
\label{rem:commutative}
The non-commutative row is real and we do not develop it. The reason is that the
system this document is anchored in has an associative-commutative parallel
composition, so its parands form a bag, not a sequence; the natural container is
$\Mset$ and the natural linear logic is the commutative one. Where genuine
sequence structure appears --- as it did when the weighted-rewriting
implementation had to add an ordered term former, and with it the corresponding
positional structural connective, to stop two distinct synapses canonicalising to
the same channel \cite{WeightedGSLT} --- the $\List$ row is what one would want,
and Lambek calculus is where it leads. That is a different paper.
\end{remark}

\section{A finer axis: which lattice operations are primitive}
\label{sec:primitive}

Principle~\ref{prin:dial} is not the end of the story, because two containers can
agree on both Booleans and still induce very different logics. The finer question
is \emph{which lattice operations the container makes primitive}, as against
merely definable or merely relational.

Consider three containers whose aggregation is union --- commutative and
idempotent in every case, so indistinguishable on the dial of
\S\ref{sec:dial} --- and yet:

\begin{itemize}[leftmargin=1.4em]
\item The opens of a \emph{topology} form a frame: arbitrary joins and finite
meets are both primitive, and the induced logic is geometric, hence
intuitionistic. A proposition is an affirmation: something that can be positively
\emph{said} \cite{Vickers89}.
\item The opens of a \emph{semitopology} \cite{Gabbay25} --- closed under
arbitrary unions but \emph{not} under intersection --- form a semiframe: only
arbitrary joins are primitive. Conjunction survives relationally and derivatively
but not as a container operation. The induced logic is about what can be
positively \emph{done}.
\item A quantale of the form $\Pfin(M)$ for a monoid $M$ has arbitrary joins
\emph{and} a multiplication inherited from $M$ --- and, if the lattice is closed
under intersection, a second, independent conjunction. Two conjunctions, from two
different levels of the container.
\end{itemize}

\begin{principle}[Primitivity]
\label{prin:primitive}
Beyond \emph{which structural rules hold} (Principle~\ref{prin:dial}), a container
is graded by \emph{which lattice operations it makes primitive}. A frame makes
arbitrary $\bigvee$ and finite $\wedge$ primitive. A semiframe makes only
$\bigvee$ primitive. A quantale makes $\bigvee$ primitive and adds a multiplication
that need not agree with any meet.
\end{principle}

\begin{figure}[h]
\centering
\begin{tikzpicture}[node distance=6mm,
  box/.style={draw,rounded corners=2pt,align=center,inner sep=5pt,minimum width=3.4cm}]
\node[box] (frame) {\textbf{frame}\\[2pt] $\bigvee$, finite $\wedge$\\[2pt]
  \emph{intuitionistic}};
\node[box,right=14mm of frame] (quant) {\textbf{quantale}\\[2pt]
  $\bigvee$, $\otimes$, ($\wedge$)\\[2pt] \emph{linear}};
\node[box,left=14mm of frame] (semi) {\textbf{semiframe}\\[2pt]
  $\bigvee$ only; $\ast$ relational\\[2pt] \emph{coalitional}};
\draw[-{Stealth}] (frame) -- node[above,font=\scriptsize]{add $\otimes$} (quant);
\draw[-{Stealth}] (frame) -- node[above,font=\scriptsize]{drop $\wedge$} (semi);
\end{tikzpicture}
\caption{Three settings of the primitivity axis. The middle is the familiar case;
the logics of interest here are on either side of it.}
\label{fig:primitivity}
\end{figure}

The rest of this part works the two ends of Figure~\ref{fig:primitivity}.

\section{Linear logic, generated}
\label{sec:linear}

\noindent\emph{In plain terms: linear logic is what you get when witnesses are
collected in a container that keeps count. Its notorious pair of conjunctions is
not a quirk of the proof theory --- it is the signature of a container with two
layers, one that multiplies and one that intersects.}

\subsection{Phase spaces are a two-layer container}

Girard's phase semantics \cite{Girard87} fixes a commutative monoid $M$ and a
distinguished subset $\bot \subseteq M$, the \emph{pole}. For $X \subseteq M$ put
\[
  X^{\bot} \;=\; \{\, y \in M : \forall x \in X.\ xy \in \bot \,\},
\]
and call $X$ a \emph{fact} when $X = X^{\bot\bot}$. The connectives are then
\[
  X \otimes Y = (X \bullet Y)^{\bot\bot}, \quad
  X \bullet Y = \{xy : x \in X,\, y \in Y\}, \qquad
  X \mathbin{\&} Y = X \cap Y, \qquad
  X \oplus Y = (X \cup Y)^{\bot\bot},
\]
with $X^{\bot}$ the linear negation, $\bindnasrepma$ and $\multimap$ recovered by duality.

Read this through \S\ref{sec:three}. The container is not one thing but a stack:
\[
  \underbrace{\Mset\,A}_{\text{inner: a monoid}}
  \quad\longrightarrow\quad
  \underbrace{\Pfin(\Mset\,A)}_{\text{outer: a lattice}}
  \quad\longrightarrow\quad
  \underbrace{\text{facts}}_{\text{the nucleus}}
\]
The inner layer is the free commutative monoid on the atoms --- a bag. It carries
a multiplication and nothing else. The outer layer is the powerset, which carries
arbitrary unions and intersections. And the pole induces a closure operator
$X \mapsto X^{\bot\bot}$ --- a nucleus --- which cuts the outer lattice down to
the facts and, in doing so, makes negation involutive.

\begin{observation}[The two conjunctions are two layers]
\label{obs:twoconj}
$\otimes$ is the inner monoid lifted pointwise and then closed; $\&$ is the outer
container's intersection. They are different connectives because they come from
different levels of the container stack, and this is why they obey different
structural rules: the inner multiplication cannot discard or duplicate a factor,
because the bag records multiplicity, while the outer intersection is set-shaped
and therefore fully structural.
\end{observation}

Observation~\ref{obs:twoconj} is the whole of the ``linear logic is generated''
claim in one line, and it does real explanatory work. The multiplicative /
additive split is presented in most introductions as a discovery about sequent
calculus --- two ways to combine contexts, one sharing and one splitting. Here it
is the difference between the two floors of a container, and the sequent-calculus
fact follows rather than being the primitive.

\begin{remark}[Quantales]
\label{rem:quantale}
Abstractly, $\Pfin(M)$ is the free quantale on the monoid $M$: a complete lattice
with an associative multiplication distributing over arbitrary joins
\cite{Rosenthal90}. A quantale equipped with a dualising element is a
\emph{Girard quantale}, and Girard quantales are exactly the algebraic models of
classical linear logic \cite{Yetter90}. Everything in this section can be run in
that generality; the phase-space presentation is the concrete case
$M = \Mset\,A$, which is the one our term structure hands us.
\end{remark}

\subsection{The separating conjunction is the multiplicative}

Now specialise the atoms to terms. In the \rhoc{} calculus parallel composition is
associative and commutative with unit $0$, so the parands of a term form a bag and
$(\T,\mid,0)$ is a commutative monoid. Lift it pointwise. What is obtained is
\[
  \sem{\varphi \mid \psi} \;=\; \{\, P \mid Q : P \real \varphi,\ Q \real \psi \,\},
\]
which is precisely the \emph{separating conjunction} of spatial logic
\cite{CairesCardelli03,MeredithRadestock05b} --- the connective that says a term
splits into a part satisfying $\varphi$ and a part satisfying $\psi$.

\begin{observation}[One connective, two names]
\label{obs:sepisotimes}
The structural connective generated by parallel composition and the multiplicative
conjunction of the induced Girard quantale are the same operation: the container's
monoid lifted pointwise. Spatial logic's $\mid$ is linear logic's $\otimes$.
Correspondingly, the \emph{graded} separation $\varphi \mid_{k}^{N} \psi$ used
throughout the mortal-scientist work --- which permits the two parts to share at
most $k$ names of a described namespace $N$ --- is a graded $\otimes$, and $k=0$
is the strict, resource-disjoint case.
\end{observation}

This is folklore-adjacent and should be positioned honestly. The logic of bunched
implications \cite{OHearnPym99} and separation logic \cite{Reynolds02} make the
same move over a \emph{partial} commutative monoid of heaps, where two heaps
compose only if their domains are disjoint, and where the additive conjunction
sits alongside the multiplicative rather than being derived from a second
container layer. The difference in our setting is instructive: the \rhoc{}
monoid is \emph{total}. Any two terms compose. So the splitting witnessing
$\varphi \mid \psi$ is not unique and not determined by a disjointness side
condition; it is chosen, and there may be many.

That non-uniqueness is not a nuisance. It is the point of contact with the
construction \cite{Zeus} is organised around. Membership in $X \bullet Y$ is an
\emph{existential over splittings} of an element, and the object of splittings of
a term is the pair $(\p\T, \T)$ --- a one-hole context together with what sits in
the hole --- which is the McBride derivative \cite{McBride01,Huet97}, the zipper,
the \emph{proto power object}. Assembling the splittings back into a collection of
the same shape as the container is the operation \cite{Zeus} calls packaging, and
its central result is that packaging is a section of the quotient presenting the
container, whose choice content is exactly the assertion that the atoms admit a
linear order.

\begin{observation}[The multiplicative is where the packaging problem lives]
\label{obs:packing}
$\otimes$ is definable relationally without any choice --- to know whether $P$
satisfies $\varphi \mid \psi$ one need only know that \emph{some} splitting works.
To produce the collection of witnessing splittings as an object of the container's
own shape requires packaging, and hence choice. The relational form is free; the
operational form is not. Linear logic's multiplicative is therefore the natural
worked instance of the proto-versus-packaged distinction, and not a separate
topic.
\end{observation}

Observation~\ref{obs:packing} will be picked up again in Part V, where the same
distinction reappears in a completely different costume: a graded where-clause
\emph{weighs} the candidate resolutions of a race without selecting one, and a
scheduler \emph{selects}. Weighing is relational and free; selecting is packaging
and is priced.

\subsection{The exponential}

The remaining connective is the exponential $!$, and on the reading above it has
an unusually concrete description: $!$ is the modality that re-licenses the
structural rules the container withheld. Contraction corresponds to idempotence
(Principle~\ref{prin:dial}); the container was chosen to be a bag precisely so as
\emph{not} to be idempotent; and $!\varphi$ marks the propositions on which one
may nonetheless duplicate. In container terms it is the passage from $\Mset$ back
up to $\Pfin$, restricted to those elements where the passage is legitimate.

\begin{remark}[What the exponential costs]
\label{rem:bangcost}
There is a price attached, and \cite{Zeus} computes it. Packaging for $\Mset$
needs a linear order on atoms, so that each finite support has a canonical sorted
representative. Packaging for $\Pfin$ needs that \emph{and} a deduplication step,
hence decidable equality on atoms. The additional equation between the two
containers is idempotence, and the additional resource is an equality test. So:
the modality that restores contraction is exactly the packaging cost of
idempotence, and that cost is a decidable equality. In a reflective calculus where
atoms are quoted processes, ``decidable equality on atoms'' is a substantive and
recognisable requirement.
\end{remark}

\begin{remark}[The exponential in the calculus, and a conjecture]
\label{rem:bangpersist}
It is tempting to identify $!$ with a syntactic feature of the calculus, and the
temptation is worth recording without being indulged. Rholang receipts come in two
flavours: linear, written \texttt{<-}, which is consumed when it fires, and
persistent, written \texttt{<=}, which is not. A persistent receipt is precisely a
capability to receive that may be used any number of times, which is the informal
gloss of $!$ on the input side; and the trophic vocabulary of the mortal-scientist
work makes the same distinction for the other side of a communication, separating
a \emph{source} --- persistent, rate-limited, non-exclusive --- from \emph{prey}
--- a linear send, exhaustible, exclusive \cite{ScientistNote}. Same sort, same
rewrite rule, different multiplicity.

There is also a suggestive piece of independent evidence. In the graded setting,
what makes a race between two messages a race at all is the \emph{linearity} of
the receipt: the loser's message survives and constitutes a which-path record, so
the two branches cannot recombine. Replace the receipt by a persistent one and
both messages are eventually consumed, the receipt survives in either branch, and
the derivations reconverge --- the interference site is destroyed
\cite{GradedWhere}. The exponential erases which-path information. That is exactly
what one would want $!$ to do to a resource-sensitive semantics, and it is
checkable rather than metaphorical.

We nevertheless leave the identification as a conjecture rather than promoting it
to a proposition. Making it precise requires a statement about contraction on
received data in the presence of reflection, which is a technical question with a
known trap in it: reflection is the mechanism by which a term can acquire
structure it did not syntactically contain, and it is already the suspect in
several open questions elsewhere in this corpus. The conjecture is that
persistence is the calculus's native $!$; the honest status is that we have the
resource reading, the cost reading, and one behavioural test, and not a proof.
\end{remark}

\subsection{Why this row matters for the rest of the document}

Two dividends, both cashed later.

First, the value dial and the collection dial turn out to be the same dial applied
at different layers. A quantale is a container of truth values; grading a formula
in a quantale, as \cite{FuzzyRho} does for the \rhoc{} calculus, means its truth
values are linear-logic truth values, with $\otimes$ for conjunction and
$\bigvee$ for disjunction. This is not a pun: it is the observation that the
algebra one uses to \emph{collect witnesses} and the algebra one uses to
\emph{value formulae} are subject to the same constraints, because in both cases
the connectives are the operations under which the algebra is closed.

Second --- and this is the payoff advertised in Part I --- the linear reading
supplies a criterion for something the graded-clause design had only by
inspection.

\begin{observation}[When the graded sort has a negation]
\label{obs:negcriterion}
The graded sort of \cite{GradedWhere} was given a semiring structure and denied
both negation and fixed points, on the grounds that $1-x$ works on $[0,1]$ but
$\Cx$ has no canonical negation and De Morgan fails. The quantale reading explains
the line. A quantale supports an involutive negation exactly when it carries a
\emph{dualising element}, i.e.\ when it is a Girard quantale
\cite{Yetter90,Rosenthal90}. Then:
\begin{itemize}[leftmargin=1.4em]
\item $[0,1]$ with \L{}ukasiewicz negation $1-x$ is a Girard quantale --- which is
why probabilistic-logic values support refutation as $1-s$, and why the mortal
forager's confirm / refute / undetermined trichotomy sits as three corners of a
square rather than requiring a separate negation in the hypothesis language;
\item a general quantale supports fixed points by Knaster--Tarski, because it is a
complete lattice --- which is why the graded fixed point is available in
quantale-graded \rhoc{} and has to be argued for separately over $\R_{\ge0}$,
where it becomes a convergence condition;
\item $\Cx$ is a rig with no order at all, hence not a quantale, hence carries
neither a dualising element nor Knaster--Tarski. Negation and fixed points are
both unavailable, and for the same reason.
\end{itemize}
So the two-sorted design is not a compromise forced by the complex case. It is the
statement that a crisp sort is needed exactly when the value algebra fails to be a
Girard quantale.
\end{observation}

Observation~\ref{obs:negcriterion} is, as far as we know, new, though it is a
short step from facts that have been available since \cite{Yetter90}. It is
recorded here rather than in Part IV because it is a fact about containers, and it
is the reason Part IV can treat negation as a single phenomenon with several
faces rather than as a recurring nuisance.

\section{The semitopological row, and the promised companion}
\label{sec:semitop}

\noindent\emph{In plain terms: if the container is the set of coalitions that can
act on their own, then two coalitions that can each act may overlap in a way that
cannot. Conjunction stops being an operation. What is left is a logic of what can
be positively done, and it is the logic of decentralised consensus.}

\subsection{The collection component}

A \emph{semitopology} \cite{Gabbay25} is a pair $(P,\Open)$ with
$\emptyset, P \in \Open$ and $\Open$ closed under arbitrary unions, but not
necessarily under intersection. Points are participants; an open is an
\emph{actionable coalition}, a set of participants with the resources to act
without anyone else's permission. A union of actionable coalitions is actionable;
an intersection need not be. That one dropped axiom is the whole difference from
topology.

Take $\Coll = \Open(P)$ and run \S\ref{sec:three}. What comes out, with nothing
added by hand, is: arbitrary disjunction is a genuine collection connective;
conjunction is \emph{not} interpretable as a container operation, though the
interior $\mathrm{int}(A \cap B)$ exists as a derived operation by union-closure;
and the primitive relational residue of conjunction is the compatibility relation
\[
  A \ast B \quad\Longleftrightarrow\quad A \cap B \neq \emptyset,
\]
recording co-realizability --- the existence of a witness realising both --- without
producing a collection of witnesses.

\begin{observation}[The framework predicts the semiframe]
\label{obs:semiframe}
From the closure properties of the container alone, the collection-component
analysis predicts that the algebraic dual of a semitopology is a complete
join-semilattice equipped not with a meet but with a symmetric, join-distributive
compatibility relation. That is exactly Gabbay's \emph{semiframe}
\cite{GabbaySemiframes}, arrived at independently as the structure needed to make
the semitopology--semiframe duality work and to express the intertwinedness
conditions governing consensus. Two routes, one relation.
\end{observation}

It is worth recording why this is the best available evidence for the whole
programme. The prediction was not tuned: the container analysis was developed for
reasons having nothing to do with distributed systems, and the structure it
predicts is idiosyncratic --- a relation where one expects an operation, with
three Horn-shaped axioms. Finding it already in the literature, derived
differently, is the kind of corroboration a generative account is supposed to
produce and a descriptive one is not.

\begin{remark}[Compatibility is a proto-conjunction]
The parallel with \S\ref{sec:linear} is exact and worth one line. The relation
$\ast$ is to conjunction what the object of splittings is to the power object: the
un-packaged, relational form that records possibility without producing the
packaged object. Recovering an actual conjunction from $\ast$ means applying an
interior --- a closure, a global operation --- just as recovering a sub-collection
from splittings meant applying $\pack$.
\end{remark}

\subsection{Points are processes: the structural and behavioural components}

\cite{Zeus} treats only the collection component of this row and defers the rest.
We supply it, following \cite{QuorumNote}.

The starting point is a puzzle in the consensus-types development
\cite{ConsensusTypes}. Its participants are channel names, an open set is a set of
channels, and the coalition modality $\exists\Box\varphi$ at a participant is
realised by a \emph{join} --- a receive that must synchronise on several channels
at once. The paper observes that the coalition modality decomposes over
\emph{participants} rather than over \emph{process structure}, and flags a
systematic comparison with spatial logic as future work.

\begin{proposition}[The distinction collapses under reflection]
\label{prop:collapse}
In the \rhoc{} calculus the two decompositions coincide. Participants are
processes; the channel names by which they are addressed are quoted processes; and
the name predicate $\quot{\varphi}$ lets a formula about process structure be
asked of a name. Consequently an open set --- a set of participants --- is a
described sub-namespace, and quantification over opens is quantification over
namespaces described by name predicates.
\end{proposition}

The proof is by inspection of the encoding and is given in \cite{QuorumNote}; what
matters here is the consequence. The structural component of the coalition logic
is namespace logic, and the behavioural component is the ordinary
context-labelled modality $\langle \Ktx_{j}\rangle$. There is no third kind of
connective. What distinguishes this row from the spatial-behavioural row is the
container, and only the container.

Writing $\Open$-indexed aggregation explicitly, the coalition modality over a
namespace $N$ and a value algebra $\V$ is
\[
  \sem{\langle E \rangle_{N}^{\V}\, \psi}
  \;=\;
  \bigoplus_{\substack{O \subseteq N \\ O \text{ open}}}\;
  \bigotimes_{q \in O} \sem{\psi}(q),
\]
--- an $\oplus$ over coalitions of an $\otimes$ within one --- and this single
display is where the three dials meet. The container supplies the quantification
pattern; the term language supplies what $\psi$ can say about each participant;
the value algebra supplies $\oplus$ and $\otimes$.

\subsection{The container forces the connective; attainment forces the algebra}

Given the display above, which $\V$ should one use? The question looks like a
modelling decision and is not.

Over $\R_{\ge0}$ with $\otimes = \times$, a coalition's value is a product of
per-participant values, so a large coalition decays toward zero: bigger quorums
would always be weaker, which inverts the intended semantics. The correct quorum
reading is that a coalition is as strong as its weakest member, i.e.\ $\otimes =
\min$ --- the tropical and Viterbi algebras.

That would be a modelling preference, except that the same two algebras are picked
out independently by a condition in Part V. A graded clause assigns a value to
each candidate; a scheduler must actually select one; and the selection is
\emph{attained} --- realised by a definable selection function returning an actual
candidate rather than an expectation --- exactly when $\oplus$ is idempotent, with
$\max$ and $\min$ the two cases of interest \cite{ChoiceFormulae}. So:

\begin{observation}[Three constraints, one algebra]
\label{obs:threeconstraints}
The container forces the shape of the coalition connective, since a meet-free
container can only offer $\oplus$-over-opens of $\otimes$-within. The intended
quorum semantics forces $\otimes = \min$. Attainment forces $\oplus$ idempotent.
The intersection of the three is the tropical/Viterbi row of
Table~\ref{tab:valuealgebras} --- which is precisely the row where a selection
function exists on the nose, so that an observer evaluating the modality receives
a \emph{witness coalition} rather than a scalar.
\end{observation}

The ecological reading of that observation is the reason it is here rather than in
an appendix. A lion pride reading which way a stampeding herd will break needs to
know not merely how strong the consensus is but \emph{which way}; a honeybee scout
committing to a nest site needs a site, not a number. The framework lands, for
independent reasons, in the one fragment where selection is definable --- which is
why an animal can act on it at all.

\subsection{What changes when the observer is mortal}

One further move separates this row from the consensus literature, and it is a
change of direction rather than of content.

The consensus-types development reads the correspondence as a term assignment:
from a proof of $\varphi$, synthesise a protocol $P$ realising it. Its open
problems are accordingly proof-theoretic --- there is no proof system for
coalition logic yet, and supplying one is most of the work. A mortal observer
needs the opposite arrow, $P \models \varphi$: given a population, has a coalition
formed? Same logic, opposite direction.

\begin{proposition}[The observation reading is proof-theory-free]
\label{prop:cheap}
Reading the coalition logic as a hypothesis language about an already-running
population requires only its semantics, which the container supplies. It requires
none of the proof-theoretic apparatus that synthesis requires.
\end{proposition}

And with that reversal comes a reframing that Part VII will take up. Gabbay's
semitopology is \emph{given}: the modeller supplies the set of actionable
coalitions. A mortal observer has to \emph{discover} it. The semitopology is
therefore not the model but the \emph{hypothesis} --- the thing the observer is
guessing at --- and coalition logic stops being a specification language and
becomes a hypothesis language. That is the sense in which a container choice is
also an epistemological commitment: it says what kind of thing the observer takes
its world to be organised into.
